\documentclass{article}
\usepackage{amsmath,amssymb}
\usepackage[vocab={solids}]{leanify}
\newtheorem{lemma}{Lemma}
\begin{document}

\section{Notation}

We write $V_n(r)$ for the volume of the ball of radius $r$ in dimension $n$.
We write $S_n(r)$ for the surface area of the sphere of radius $r$ in dimension $n$.
We write $C(r,h)$ for the volume of the cylinder with radius $r$ and height $h$.
We write $K(r,h)$ for the volume of the cone with radius $r$ and height $h$.
We write $T(R,r)$ for the volume of the torus with radii $R$ and $r$.
We write $Q(a)$ for the volume of the regular octahedron with edge $a$.
We write $\Sigma_n$ for the volume of the standard simplex in dimension $n$.

\section{Balls}

\begin{lemma}\label{balls}
$V_1(1)=2$, $V_2(1)=\pi$, $V_3(1)=\frac{4\pi}{3}$ and $S_3(1)=4\pi$.
\end{lemma}

\begin{lemma}\label{derivative}
Let $r>0$ be a real number. Then $S_2(r)=2\pi r$.
\end{lemma}

\section{Solids}

\begin{lemma}\label{solids}
$C(1,3)=3\pi$, $K(1,3)=\pi$, $T(2,1)=4\pi^2$, $Q(1)=\frac{\sqrt2}{3}$ and $\Sigma_3=\frac16$.
\end{lemma}

\end{document}
